\documentclass[11pt]{article}
\usepackage{fontspec}
\usepackage[french]{babel}
\usepackage{amsmath}
\usepackage{amssymb}
\usepackage{amsthm}
\usepackage{mathrsfs}
\usepackage{tikz-cd}
\usepackage{enumitem}

\newcommand{\printedpage}[1]{\marginpar{\footnotesize p.~#1}}
\newcommand{\msp}[1]{\marginpar{\footnotesize [#1]}}
\newcommand{\Hom}{\operatorname{Hom}}
\newcommand{\Ob}{\operatorname{Ob}}
\newcommand{\pr}{\operatorname{pr}}
\newcommand{\card}{\operatorname{card}}
\newcommand{\Fl}{\operatorname{Fl}}
\newcommand{\Filt}{\operatorname{Filt}}
\newcommand{\Ind}{\operatorname{Ind}}

\begin{document}
\setcounter{footnote}{6}

\printedpage{9}
% head: A. Grothendieck : Dérivateurs, Ch. XVIII, Édition des Universités de Montpellier II et Paris VII

\noindent\textit{De plus, si c est satisfait, on a ceci : pour $s\in S$, $t\in T$, considérons}
\[
(s,t)\ \in\ (S\times T)\ \subset\ (S\times T)^\wedge\ \simeq\ S^\wedge\boxtimes T^\wedge\ ,
\]
\textit{qui s'explicite aussi comme}
\[
\left\{
\begin{aligned}
&\rlap{$\displaystyle (s,t)\ =\ \pr_1^*(s)\times\pr_2^*(t)\ \simeq\ s\overset{S^\wedge,T^\wedge}{\boxtimes}t\in\Ob\,(S\times T)^\wedge\ ,$}\\
&\textit{i.e.}\ (s,t)(s',t')\ =\ s(s')\times t(t') &&=\ \Hom_S(s',s)\times\Hom_T(t',t)\\
& &&=\ \Hom_{S\times T}((s',t'),(s,t))\ ,
\end{aligned}
\right.
\tag*{(4.3.1.4)}
\]
\[
\begin{tikzcd}
S\times T \arrow[r, "\mathrm{pr}_2"] \arrow[d, "\mathrm{pr}_1"] & T\\
S\ ,
\end{tikzcd}
\]
\msp{page 66}\textit{on a}
% note: the label (4.3.1.4) is printed for this display as well as for the preceding one
\[
s\overset{\mathcal M,\mathcal N}{\boxtimes}t\ \simeq\ \varphi(s,t)\ ,
\tag*{(4.3.1.4)}
\]
\textit{et $\varphi(S\times T)=\{s\boxtimes t\mid s\in S,\ t\in T\}$ est une partie strictement génératrice de $\mathcal M\boxtimes\mathcal N$ (}\footnote{\textbf{N.B.} On a bien sûr $(s,t)=s\overset{S^\wedge,T^\wedge}{\boxtimes}t$.}\textit{).}

\noindent\textbf{N.B.} Je ne connais pas d'exemple où les conditions envisagées ne soient satisfaites, mais doute pourtant à présent que même la plus faible le soit toujours, i.e. que $\mathcal M\boxtimes\mathcal N$ soit un pseudo-topos dès que $\mathcal M$ et $\mathcal N$ le sont.

\noindent\textsc{Démonstration de} 4.3.1. L'implication b $\Longrightarrow$ b$'$ est tautologique. L'implication b$'$ $\Longleftrightarrow$ c résulte du

% marginal vertical bar printed beside Lemme 4.3.2 (from "Lemme 4.3.2." to "soit représentable.")
\noindent\textbf{Lemme 4.3.2.} \textit{Soient $\mathcal E$, $X$ des $\mathfrak U$-catégories avec $X$ petite, $\mathcal E$ stable par petites $\varprojlim$ ,}
\[
\alpha\ :\ \mathcal E\longrightarrow X^\wedge
\]
\textit{un foncteur. Pour que $\alpha$ admette un adjoint à gauche $\varphi$, il faut et il suffit que pour tout $x\in\Ob X$, le foncteur}
\[
\left\{
\begin{array}{ccl}
\xi & \longmapsto & \Hom_{X^\wedge}(x,\alpha(\xi))\ =\ \alpha(\xi)(x)\\
\mathcal E^\circ & \longrightarrow & \mathrm{Ens}
\end{array}
\right.
\]
\textit{soit représentable.}

\noindent C'est évidemment nécessaire. La suffisance résulte du fait que $X\subset X^\wedge$ est stricte%
% foot: Version 17 décembre 2025
\printedpage{10}%
% head: A. Grothendieck : Dérivateurs, Ch. XVIII, Édition des Universités de Montpellier II et Paris VII
\msp{page 67}%
ment génératrice dans $X^\wedge$, i.e. que tout objet $F$ de $X^\wedge$ est limite inductive $F=\varinjlim\limits_I F_i$, les $F_i$ dans $X$ (on prend $I=X/F$), d'où
\[
\begin{aligned}
\Hom(F,\alpha(\xi))\ &\simeq\ \varprojlim_I\ \underbrace{\Hom(F_i,\alpha(\xi))}_{\Hom(\rho(F_i),\xi)}\\
&\simeq\ \Hom(\varinjlim_I\,(\rho(F_i),\xi))\ ,
\end{aligned}
\]
donc le foncteur $\xi\longmapsto\Hom(F,\alpha(\xi))$ est représenté par $\rho(F)\overset{\text{déf}}{=}\varinjlim\limits_I\rho(F_i)$.

\noindent Prouvons c $\Longrightarrow$ a. Pour l'existence des petites $\varinjlim$ dans $\mathcal M\boxtimes\mathcal N$, et l'existence d'une partie strictement génératrice, on invoque le lemme général :

% marginal vertical bar printed beside Lemme 4.3.3 (from "Lemme 4.3.3." to "comme")
\noindent\textbf{Lemme 4.3.3.} \textit{Soit $\alpha:\mathcal E\hookrightarrow\mathcal F$ un foncteur pleinement fidèle admettant un adjoint à gauche $\rho$.}
\begin{enumerate}[label={}]
\item[\textcircled{a}] \textit{Si les $\varinjlim$ d'un type donné $I$ existent dans $\mathcal F$, elles existent dans $\mathcal E$, et se calculent comme}
\[
\varinjlim_I{}_{\mathcal E}\,\xi_i\ =\ \rho(\varinjlim_I{}_{\mathcal F}\,\alpha(\xi_i))\ .
\tag*{(4.3.3.1)}
\]
\item[\textcircled{b}] \textit{Si $U$ est une petite partie strictement génératrice de $\mathcal F$, $\rho(U)$ est une petite partie strictement génératrice dans $\mathcal F$.}
\end{enumerate}

\noindent En effet, si $\xi$ désigne le deuxième membre \texttt{[de l'égalité 4.3.3.1]}, et $\eta$ un objet variable de $\mathcal E$, on aura
\[
\begin{aligned}
\Hom_{\mathcal E}(\xi,\eta)\ &\overset{\text{adj.}}{\simeq}\ \Hom_{\mathcal F}(\varinjlim_I\alpha(\xi_i),\alpha(\eta))\\
&\simeq\ \varprojlim_I\quad\underbrace{\Hom_{\mathcal F}(\alpha(\xi_i),\alpha(\eta))}_{\simeq\,\Hom_{\mathcal E}(\xi_i,\eta)\ \ \text{car $\alpha$ pleinement fidèle}}\quad .
\end{aligned}
\]
\msp{page 68}
\[
\simeq\ \ \varprojlim_I\Hom_{\mathcal E}(\xi_i,\eta)\ ,
\]
ce qui prouve que $\xi$ est une limite inductive des $\xi_i$ dans $\mathcal E$ (\footnote{Pour prouver b), on note que la pleine fidélité de $\alpha$ implique que $\rho$ est essentiellement surjectif ; comme $\rho$ commute aux petites $\varinjlim$ et que tout objet de $\mathcal F$ est de la forme $\varprojlim\limits_I x_i$, les $x_i$ dans $U$, on aura $\rho(x)=\varinjlim\limits_I\rho(x_i)$, les $\rho(x_i)$ dans $\rho(U)$, on gagne.}).

\noindent Ceci prouve que (moyennant c)) $\mathcal M\boxtimes\mathcal N$ est stable par petites $\varinjlim$ , donc il reste à prouver que $\mathcal M\boxtimes\mathcal N$ admet une petite famille strictement génératrice. Cela résultera du b) dans le lemme, puisque $S\times T$ engendre strictement $(S\times T)^\wedge$. Cela prouve en même temps la dernière assertion de la proposition, sous réserve de prouver la formule (4.1.3.4). Il faut donc prouver un isomorphisme fonctoriel en $F\in\Ob\mathcal M\boxtimes\mathcal N$, % [sic] (4.1.3.4) -> (4.3.1.4), the formula s ⊠ t ≃ φ(s,t) to be proved is labelled (4.3.1.4) on p. 9
\[
\Hom(\rho(s,t),F)\ \simeq\ F(s,t)\ ,
\]
mais par adjonction on a
\[
\Hom(\rho(s,t),F)\ \simeq\ \Hom((s,t),\alpha(F))\ =\ \alpha(F)(s,t)\ =\ F(s,t)\ ,
\]
OK.

% foot: Version 17 décembre 2025
\printedpage{11}
% head: A. Grothendieck : Dérivateurs, Ch. XVIII, Édition des Universités de Montpellier II et Paris VII

\noindent\msp{page 69}\textbf{Proposition 4.3.4.} \textit{Soit $\mathcal M$ une $\mathfrak U$-catégorie. Les conditions suivantes sont équivalentes.}
\begin{enumerate}[label={}]
\item[\textit{a)}] \textit{$\mathcal M$ est une catégorie $\mathfrak U$-pseudo-topos, i.e. elle est stable par $\mathfrak U$-limites inductives, et admet une $\mathfrak U$-petite sous-catégorie $S$ strictement génératrice.}
\item[\textit{b)}] \textit{Il existe une $\mathfrak U$-petite catégorie $S$, et un foncteur $i:\mathcal M\hookrightarrow S^\wedge$} pleinement fidèle, \textit{ayant un adjoint à gauche.}
\end{enumerate}

\noindent\textsc{Démonstration}. On a a) $\Longrightarrow$ b), car il suffit de prendre pour $S$ dans b) celui spécifié dans a), l'hypothèse sur $S$ implique que le foncteur canonique
\[
\mathcal M\ \longrightarrow\ S^\wedge\ ,\qquad\quad \text{composé de } \mathcal M\hookrightarrow\ \mathcal M^\wedge\ \xrightarrow{\ \text{restriction}\ }S^\wedge
\]
est pleinement fidèle. Qu'il admette un adjoint à gauche résulte du lemme 4.3.2. L'implication b) $\Longrightarrow$ a) résulte du lemme 4.3.3.

\noindent\textbf{Corollaire 4.3.5.} \textit{Soit $\mathcal M$ une $\mathfrak U$-catégorie. Pour que $\mathcal M$ soit un $\mathfrak U$-pseudo-topos} accessible \textit{(i.e. dont tous les objets sont accessibles), il faut et il suffit qu'il existe $S$ dans} Cat \textit{et un foncteur pleinement fidèle et} accessible $\mathcal M\overset{\alpha}{\hookrightarrow}S^\wedge$\textit{, admettant un adjoint à gauche $\rho$.}

\noindent\msp{page 70}\textsc{Démonstration}. Il suffit de noter que si $S$, $\alpha$ sont comme dans 4.3.4 b), alors $\alpha$ est accessible si et seulement si les $\rho(S)$ sont formées d'objets accessibles dans $\mathcal M$ ; comme cette partie est strictement génératrice, on sait que cela implique que \textit{tout} objet de $\mathcal M$ est accessible. (Cf. SGA 4 I 9.9.)

\noindent Je vais essayer de donner un énoncé plus complet :

\noindent\textbf{Théorème 4.3.6.} \textit{Soit $\mathcal M$ une $\mathfrak U$-catégorie. Conditions équivalentes :}
\begin{enumerate}[label={}]
\item[\textit{a)}] \textit{$\mathcal M$ stable par petites $\varinjlim$ , et elle admet une petite sous-catégorie pleine $S$, strictement génératrice, et formée d'objets accessibles de $\mathcal M$.}
\item[\textit{a$'$)}] \textit{$\mathcal M$ est une catégorie $\mathfrak U$-pseudo-topos faiblement accessible, i.e. satisfait a), avec en plus} tout \textit{élément de $\mathcal M$ (pas seulement ceux de $S$) accessible dans $\mathcal M$.}
\item[\textit{b)}] \textit{Il existe une petite catégorie $S$, et un foncteur $\alpha:\mathcal M\hookrightarrow S^\wedge$} pleinement fidèle et accessible, \textit{tel que $\alpha$ admette un adjoint à gauche.}
\item[\textit{c)}] \msp{page 71}\textit{$\mathcal M$ est stable par petites $\varinjlim$ , est accessible, et admet une filtration cardinale} $(\Filt^\pi\mathcal M)_{\pi\geqslant\pi_0}$.
\item[\textit{d)}] \textit{Il existe un cardinal $\pi_1$, une petite catégorie $\mathcal M_{\pi_1}$, stable par limites inductives de cardinal $\leqslant\pi_1$, tels qu'on ait une équivalence de catégories}
\[
\mathcal M\ \xleftarrow{\ \sim\ }\ \Ind_{\pi_1}\mathcal M_{\pi_1}\ ,
\tag*{(4.3.6.1)}
\]
\textit{où le deuxième membre désigne la sous-catégorie pleine de} $\Ind(\mathcal M_{\pi_1})$ \textit{formée des ind-objets $(X_i)_{i\in I}$ de $\mathcal M_{\pi_1}$ tels que $I$ soit grand devant $\pi_1$.}
\end{enumerate}

\noindent\textsc{Démonstration de} 4.3.6 (\footnote{Pour une démonstration complète (celle donnée ici canule à partir d'un certain moment), cf. §4.6 (p. 139-148).}). On a déjà vu a$'$ $\Longrightarrow$ a $\Longrightarrow$ b $\Longrightarrow$ a$'$, i.e. l'équivalence de a, a$'$, b, soit PTA cette condition. Nous allons prouver les implications
\[
\mathrm{PTA}\ \Longrightarrow\ \mathrm{c}\ \Longrightarrow\ \mathrm{d}\ \Longrightarrow\ \mathrm{PTA}\ ,
\]
% foot: Version 17 décembre 2025
\printedpage{12}%
% head: A. Grothendieck : Dérivateurs, Ch. XVIII, Édition des Universités de Montpellier II et Paris VII
ce qui établira le théorème.

\noindent\textit{Prouvons PTA $\Longrightarrow$ c.} Soit $S$ une petite sous-catégorie pleine strictement génératrice. Ainsi, \msp{page 72}pour tout $x$ dans $\mathcal M$ on a
\[
x\ =\ \varinjlim_{s\in S/x}s
\tag*{(4.3.6.2)}
\]
(cf. SGA 4 I 7.2 (i)). Soit $\pi_0=\card\Fl(S)$, et pour tout cardinal $\pi\geqslant\pi_0$, soit
\[
\begin{array}{lcl}
\Filt^\pi\mathcal M & = & \text{sous-catégorie pleine de $\mathcal M$ formée des objets $x$ tels que}\\
 & & \card\,(\Fl(S/x))\leqslant\pi\ .
\end{array}
\tag*{(4.3.6.3)}
\]
Je veux prouver que l'on trouve ainsi une filtration cardinale sur $\mathcal M$, i.e. que les conditions a) b) c) de SGA 4 I 9.12 sont satisfaites. C'est prouvé dans loc. cit. 9.13 sous des hypothèses légèrement différentes, la condition d'accessibilité de $\mathcal M$ étant remplacée par celle (beaucoup moins sympathique!) que les familles de flèches épimorphiques strictes dans $\mathcal M$ sont épimorphiques strictes universelles. Cette condition n'est utilisée que dans la démonstration de la condition b), à savoir la stabilité de $\Filt^\pi\mathcal M$ par $\varinjlim\limits_I$ de cardinal $\leqslant\pi$ (sans avoir d'ailleurs à supposer $I$ filtrante, ni a fortiori grand devant $\pi_0$). Il faut d'ailleurs dire que la définition de $\Filt^\pi\mathcal M$ donnée dans loc. cit. n'est pas tout à fait la même que (4.3.6.3), \msp{page 73}mais c'est
% the printed line break "car-/dinal" inside the display is kept as printed
\[
\begin{array}{lcl}
\Filt'^{\pi}\mathcal M & = & \text{sous-catégorie pleine de $\mathcal M$ formée des objets $x$ tels qu'il}\\
 & & \text{existe une famille épimorphique stricte $x_i\xrightarrow{\alpha_i}x$ de car-}\\
 & & \text{dinal $\leqslant\pi$, avec $x_i\in S\ \forall\ i\in I$.}
\end{array}
\tag*{(4.3.6.3$'$)}
\]
Ainsi on a
\[
\Filt^\pi\mathcal M\ \subset\ \Filt'^{\pi}\mathcal M\ ,
\]
mais dans loc. cit. on prouve que pour $x$ dans $\Filt'^{\pi}\mathcal M$, on a
\[
\card\Fl(S/x)\ \leqslant\ \pi^{\pi_0}\ ,
\]
d'où
\[
\Filt^\pi\mathcal M\ \subset\ \Filt'^{\pi}\mathcal M\ \subset\ \Filt^{\pi^{\pi_0}}\mathcal M\ ,
\]
donc si $\pi=\pi^{\pi_0}$, p.ex. $\pi=2^c$ avec $c\geqslant\pi_0$, alors on trouve
\[
\Filt^\pi\mathcal M\ =\ \Filt'^{\pi}\mathcal M\ .
\]

\noindent Donc les deux filtrations on l'air $\pm$ équivalentes. Mais la démonstration donnée de la stabilité de $\Filt'^{\pi}$ semble vraiment faire appel de façon essentielle à la condition d'universalité des familles épimorphiques strictes (donc épimorphiques effectives, avec les hypothèses faites impliquant l'existence de petites $\varprojlim$ ). Donc il me faut bon gré, mal gré, refaire le travail! % [sic] on l'air -> ont l'air, verb avoir (3rd person plural) intended

% foot: Version 17 décembre 2025
\printedpage{13}
% head: A. Grothendieck : Dérivateurs, Ch. XVIII, Édition des Universités de Montpellier II et Paris VII

\noindent\msp{page 74}Condition a) : $\Filt^\pi$ équivalente à une petite catégorie. C'est immédiat, puisque par (4.3.6.2), $\Filt^\pi$ est formée de $\varprojlim\limits_I$ dans $\mathcal M$ d'objets de $S$ (petite catégorie), avec $I$ telle que $\card\Fl I\leqslant\pi$.

\noindent Condition b) : $\Filt^\pi$ stable par $\varprojlim\limits_I$ , quand $\card\Fl I\leqslant\pi$ (\footnote{quand $I$ filtrante et grande devant $\pi_0$ -- je ne le prouve que dans ce cas, hélas!}). Soit donc $(x_i)_{i\in I}$, où
\[
\begin{array}{ccc}
I & \longrightarrow & \mathcal M\\[1ex]
i & \longmapsto & x_i
\end{array}
\]
un système inductif dans $\mathcal M$, avec des $x_i$ dans $\Filt^\pi$. Considérons le produit fibré
\[
\begin{tikzcd}
I \arrow[d, "\varphi"'] & S/I = J \arrow[l, "\text{cofibré}"'] \arrow[d] \arrow[dr, dotted, "\psi"] & & \\
\mathcal M & S/\mathcal M \arrow[l, "\text{but}"', "\text{Cat-cofibré}"] \arrow[r, "\text{source}"] & S \arrow[r, hook] & \mathcal M\ ,
\end{tikzcd}
\tag*{(4.3.6.4)}
\]
où $J$ est Cat-cofibré sur $I$, sa fibre en $i\in\Ob I$ étant isomorphe à $S/x_i$, donc $\card\Fl J_i\leqslant\pi$. Comme par hypothèse $\card\Fl I\leqslant\pi$, on en conclut que $\card\Fl J\leqslant\pi^3=\pi$. (\textbf{N.B.} L'ensemble des flèches $v:j_0\longrightarrow j_1$ de $J$ au dessus d'une flèche $u:i_0\longrightarrow i_1$ de $I$ est “contenu” dans l'ensemble produit $\Ob J_{i_0}\times\Fl J_{i_1}$, par l'application
\[
v\longmapsto(\operatorname{source}v,v':u_*(j_0)\longrightarrow j_1)\ ,
\]
donc de cardinal majoré par $\pi^2$.) On \msp{page 75}a alors
\[
x\ \overset{\text{déf}}{=}\ \varinjlim_I x_i\ =\ \varinjlim_J\ \underbrace{\xi_j}_{=\,\psi(j)}\qquad\quad\mbox{(\footnotemark)}\ ,
\]
\footnotetext{On utilise la propriété de distributivité des $\varinjlim$ , relative à un foncteur Cat-cofibrant
\[
\begin{tikzcd}[ampersand replacement=\&]
J \arrow[r] \arrow[dr] \& I\\
\& \mathcal M\ ,
\end{tikzcd}
\]
explicitée dessus.}%
où cette fois la limite est une limite d'objets de $S$, prise dans $\mathcal M$. Il s'ensuit que $x\in\Filt'^{\pi}$, ce n'est pas tout à fait $x\in\Filt^\pi$. (Ça implique $x\in\Filt^{\pi^{\pi_0}}$, donc $x\in\Filt^\pi$ si $\pi=\pi^{\pi_0}$, p.ex. $\pi=2^c$ avec $c\geqslant\pi_0$.) Mais c'est le moment d'utiliser le fait que les objets de $S$ sont accessibles, et on prendra $\pi_0$ tel que les objets de $S$ soient $\pi_0$-accessibles, et cette fois on prendra quand même $I$ filtrante et grande devant $\pi_0$ (pour nous borner à ce qui est exigé dans la condition b) de loc. cit.). On doit prouver que $\card\,(\Fl(S/x))\leqslant\pi$ (ce qui signifie $x\in\Filt^\pi$). Mais comme les objets de $S$ sont $\pi_0$-accessibles et $I$ grande devant $\pi_0$, on voit aussitôt que
\[
S/x\ \simeq\ \varinjlim_i S/x_i\ ,
\]
% foot: Version 17 décembre 2025

\end{document}
